% problem-000074 8 有机反应机理
\section*{8. 有机反应机理}
\textit{下列为分问。}

\subsection*{8-1}
现有如下待完成的反应：

$
\begin{array}{c} \mathrm{D} _ {\mathrm{H}} \\ \mathrm{H} _ {3} \mathrm{C} \end{array} \xrightarrow [ \text {Br} ]{\text {Br} _ {2}} \xrightarrow [ h \nu ]{}
$

\subsection*{8-1}
-1 判断该反应物⼿性中⼼的构型。

\subsection*{8-1}
-2 画出此单溴化反应所有产物的结构简 $\tt { \overline { { I } } } \hat { \zeta } _ { \mathrm { < } }$

\subsection*{8-2}
变⾊眼镜⽚在阳光下显深⾊，从⽽保护⼈的眼睛；当⼈进⼊室内后，镜⽚逐步转变为⽆⾊透明，有利于在弱光下观察周边事物。⽬前流⾏的⼀种变⾊眼镜的变⾊原理如下：

（图：）

判断上述变⾊过程中反应条件c1和 ${ \mathrm { c 2 } } .$ ，简述理由。

\subsection*{8-3}
维⽣素C普遍存在于动植物中，在动物肝脏中以葡萄糖为原料经四步反应合成；⼯业⽣产的⽅式也是以葡萄糖经以下过程合成的：

（图：）

\subsection*{8-3}
-1 写出反应a和b的具体条件（可以不写溶剂）。

\subsection*{8-3}
-2 画出 L-⼭梨糖和化合物 C 的 Fischer 投影式。

\subsection*{8-4}
以下给出四个取代反应，右图为其中某⼀反应的反应势能图：

\subsection*{(1)}
$\mathrm { C H _ { 3 } C H _ { 2 } B r + N a O C H _ { 3 }  }$

\subsection*{(2)}
$\mathrm { ( C H _ { 3 } ) _ { 3 } C B r + H O C H _ { 3 }  }$

\subsection*{(3)}
$( \mathrm { C H _ { 3 } } ) _ { 2 } \mathrm { C H I } + \mathrm { K B r } $

\subsection*{(4)}
$\mathrm { ( C H _ { 3 } ) _ { 3 } C C l + ( C _ { 6 } H _ { 5 } ) _ { 3 } P  }$

（图：）

\subsection*{8-4}
-1 指出哪个反应与此反应势能图相符。

\subsection*{8-4}
-2 画出过渡态E和F的结构式。
